\section{Apéndice de terminología y método: explicación de ingeniería de los conceptos frecuentes en la entrevista}

Este apéndice unifica los términos frecuentes de la entrevista en el formato «definición + malentendido común + punto de verificación de implementación», para que los equipos compartan el mismo significado internamente y se reduzca el costo de comunicación entre equipos. Muchos proyectos avanzan lentamente no porque falte tecnología, sino porque la inconsistencia terminológica provoca un desalineamiento de objetivos.

\begin{longtable}{p{0.16\textwidth}p{0.30\textwidth}p{0.46\textwidth}}
\toprule
Término & Definición concisa & Punto de verificación de ingeniería (para evitar quedarse en eslogan) \\
\midrule
\endfirsthead
\toprule
Término & Definición concisa & Punto de verificación de ingeniería (para evitar quedarse en eslogan) \\
\midrule
\endhead
Preentrenamiento & Entrenamiento de predicción del siguiente token sobre un corpus general a gran escala & ¿Existe una estrategia de equilibrio de corpus entre dominios? ¿Se registran el origen y la versión del corpus? ¿Se monitorea la degradación de las capacidades básicas? \\
Mid-training & Etapa de moldeo de capacidades específicas entre el preentrenamiento y el posentrenamiento & ¿Se construyen datos y objetivos dirigidos para contexto largo, uso de herramientas o dominios lingüísticos específicos? \\
Post-training & Etapa de posprocesamiento que incluye alineación de instrucciones, optimización de preferencias y aprendizaje por refuerzo, entre otros & ¿Se separan los objetivos de «aumento de capacidad» y «alineación de estilo» para evitar que una sola métrica induzca a error? \\
RLHF & Aprendizaje de una función de recompensa a partir de preferencias humanas y optimización de la política & ¿Se evalúa el sesgo del modelo de recompensa? ¿Existe monitoreo de reward hacking? \\
RLVR & Aprendizaje por refuerzo con recompensas de corrección verificable & ¿Es estable el verificador? ¿La cobertura de tareas es suficiente? ¿Se evita la contaminación entre entrenamiento y evaluación? \\
Inference-time scaling & Aumentar el presupuesto de cómputo durante la inferencia para mejorar la calidad & ¿Existe una estrategia de enrutamiento clara? ¿La ganancia de calidad compensa el aumento de latencia y costo? \\
MoE & Arquitectura dispersa que activa mediante enrutamiento un pequeño número de subredes expertas & ¿Realmente se ahorra costo de inferencia por unidad? ¿Son controlables la estabilidad del enrutamiento y el equilibrio de carga? \\
MLA/GQA & Variantes del mecanismo de atención que optimizan la caché y las consultas & ¿Aumenta el rendimiento con contexto largo? ¿Se mantiene la precisión en tareas clave? \\
KV cache & Almacenamiento en caché de las claves y los valores de atención durante la inferencia para reducir el cómputo repetido & ¿Son observables la tasa de aciertos de la caché, el uso de memoria y la estrategia de reutilización entre solicitudes? \\
Context compaction & Compresión del contexto largo en un resumen breve para usarlo en la inferencia posterior & ¿La tasa de recuperación de información clave tras la compresión alcanza el estándar? ¿Se admite regresar al texto original? \\
Tool Use & El modelo invoca herramientas externas de búsqueda, cálculo o ejecución & ¿Se minimizan los límites de permisos? ¿Los registros de invocación de herramientas admiten auditoría? \\
Agentic coding & Flujo de un agente de código que planifica, ejecuta y corrige en varios pasos & ¿Existe una plantilla de especificación, una aceptación automática y un ciclo cerrado de reversión ante fallos? \\
Spec-driven development & Se define primero una especificación ejecutable y luego se impulsa la implementación a partir de ella & ¿La especificación es verificable, aceptable y versionable? ¿Cubre los flujos de excepción? \\
Data mixture & Estrategia de entrenamiento que mezcla distintas fuentes de corpus en determinadas proporciones & ¿La mezcla se ajusta dinámicamente con base en los resultados de la evaluación? ¿Existen experimentos previos con modelos pequeños? \\
Synthetic data & Muestras de entrenamiento generadas por un modelo o un programa & ¿Se filtran mediante verificación? ¿Se registran la cadena de generación y los indicadores de calidad? \\
Data pollution & Los datos de entrenamiento se contaminan con contenido generado por IA de baja calidad & ¿Se monitorean la tasa de duplicación, la consistencia y la deriva en la distribución de las fuentes? \\
Benchmark overfitting & Optimizar para la tabla de clasificación en lugar de mejorar la capacidad real & ¿Existe una doble verificación mediante evaluaciones ocultas y tareas reales en producción? \\
Jagginess & La capacidad del modelo es muy desigual entre tareas & ¿Se evalúa por clústeres de tareas en lugar de con una sola puntuación global? ¿Se publica el espectro de fallos? \\
Model router & Planificador que selecciona el modelo o el modo de inferencia según la tarea & ¿Las reglas de enrutamiento son interpretables? ¿La estrategia de degradación es explícita? \\
Latency budget & Presupuesto de latencia aceptable para el usuario & ¿Se establecen umbrales distintos según el tipo de tarea? ¿Se ofrece retroalimentación predecible? \\
Cost-to-serve & Costo de servicio de extremo a extremo por solicitud & ¿Se desglosa el costo por capa funcional? ¿Existen alertas de costo y un mecanismo de retroalimentación de políticas? \\
Safety policy & Restricciones de salida y mecanismos de manejo para escenarios de alto riesgo & ¿Existen reglas de línea roja, un proceso de escalamiento y un mecanismo de intervención humana? \\
Model spec & Documento con las normas de comportamiento que debe seguir el modelo & ¿Lo referencian en conjunto el entrenamiento, la evaluación y el proceso de lanzamiento, en lugar de ser un documento aislado? \\
Open-weight model & Forma de publicación de modelos con pesos descargables e implementables localmente & ¿Los términos de licencia son claros? ¿Las empresas pueden usarlo comercialmente y adaptarlo de forma legal? \\
API lock-in & Dependencia de la ruta hacia la interfaz de un único proveedor & ¿Se practican simulacros con modelos alternativos? ¿La capa de abstracción de la interfaz es migrable? \\
Observabilidad & Capacidad de medir y diagnosticar el comportamiento del modelo y del sistema & ¿Cubre toda la cadena: entrada, enrutamiento, salida, invocación de herramientas y repetición de fallos? \\
Rollback & Capacidad de volver a una versión estable cuando ocurre un riesgo en producción & ¿La reversión está disponible con un solo clic? ¿Las versiones históricas se mantienen en un estado funcional? \\
Human-in-the-loop & Participación humana en decisiones clave o etapas de revisión & ¿Está claro quién interviene y cuándo? ¿Se cuantifican el beneficio y el costo de la intervención? \\
Agency & Capacidad activa de las personas para controlar sus objetivos y acciones & ¿El diseño del producto conserva el poder de decisión del usuario, en lugar de que este acepte pasivamente las acciones del modelo? \\
Governance-by-design & Incorporar los requisitos de gobernanza desde la etapa de diseño del sistema & ¿La auditoría de cumplimiento forma parte del proceso de desarrollo, en lugar de ser un parche posterior al lanzamiento? \\
\bottomrule
\end{longtable}

\begin{knowledgebox}{Por qué el apéndice también es importante}
En la colaboración a gran escala, la falta de uniformidad terminológica hace que una misma palabra represente objetivos distintos en cada equipo, lo que termina generando conflictos entre las métricas de evaluación, lanzamiento y operación. Convertir la terminología en ingeniería es la condición previa para que un sistema complejo pueda volverse colaborativo.
\end{knowledgebox}

\subsection{Resumen de la sección}

Unificar la terminología no es un trabajo documental, sino ingeniería de sistemas. Solo cuando los conceptos se convierten en puntos de verificación, el equipo puede mantener una dirección coherente durante una iteración a alta velocidad.

